% TOP_PROBLEM: {"id":30007653,"problem_number":"LOCAL-30007653","title":"Scaling durable graduation programs","category_id":21,"category":{"id":21,"name":"economics","display_name":"Economics","description":"Open economic theory statements, published research agendas, and proposed scoped specifications; question provenance and historical priority are qualified per record.","slug":"economics","order_index":21,"created_at":"2026-10-08T00:00:00Z"},"status":"research_question","external_url":null,"legacy_ids":[],"published":false,"statement":"In the explicitly specified setting: Which parts of multifaceted graduation programs are essential for durable and intergenerational poverty reduction when governments deliver them at national scale? The mathematical statement above fixes the target, assumptions and scope of this catalogue formulation.","economics":{"original_id":142,"field":"Development and poverty","kind":"design","formalization_basis":"adapted_research_question","source_status":"explicit_open","priority_status":"earliest_found","priority_note":"Earliest directly inspected qualifying passage in this audit; earlier origins were not established. A review or research-agenda statement does not imply original priority.","earliest_verified_reference":{"authors":["Innovations for Poverty Action"],"title":"The Ultra-Poor Graduation Approach: Research and Learning Agenda for the Next Wave of Graduation Programs","year":2025,"url":"https://poverty-action.org/graduation-approach","doi":null,"locator":"Section \"What’s Next for Graduation Evidence\", paragraphs 2–5","evidence_summary":"The agenda explicitly asks which components are necessary, which adaptations reduce costs, and how government delivery scales."},"current_reference":{"authors":["Innovations for Poverty Action"],"title":"The Ultra-Poor Graduation Approach: Research and Learning Agenda for the Next Wave of Graduation Programs","year":2025,"url":"https://poverty-action.org/graduation-approach","doi":null,"locator":"Section \"What’s Next for Graduation Evidence\", paragraphs 2–5","evidence_summary":"The agenda explicitly asks which components are necessary, which adaptations reduce costs, and how government delivery scales."},"formalization_note":"The cited passage establishes a research direction, not this exact mathematical statement. The notation, population, policy menu, assumptions, and estimands are an original adaptation for this catalogue; no universal unsolved theorem or source priority is claimed."},"statusQualification":"Published question or research agenda verified at the cited locator. The mathematical specification is adapted for this catalogue and includes additional assumptions; source publication does not establish first-ever priority or certify this exact model remains unsolved. The cited passage establishes a research direction, not this exact mathematical statement. The notation, population, policy menu, assumptions, and estimands are an original adaptation for this catalogue; no universal unsolved theorem or source priority is claimed.","statusReviewedAt":"2026-10-08"}
% FIRST_POSTED: 2026-10-08
% ENTRY_KIND: statement_only
% !TeX program = lualatex
\documentclass[11pt]{article}
\usepackage{fontspec}
\usepackage[a4paper,margin=25mm]{geometry}
\usepackage{amsmath,amssymb,mathtools}
\usepackage{xurl}
\usepackage[hidelinks]{hyperref}
\newcommand{\sref}[2]{\hyperlink{source:#1:#2}{[S#2]}}
\setlength{\emergencystretch}{3em}
\begin{document}
% problemId:30007653
\section*{Scaling durable graduation programs}\label{30007653}
\subsection{Definitions and mathematical statement}
Fix a target population and a government delivery system for an economic-inclusion program. Let \(p=(a,t,c,s)\) specify the quantities of productive assets, training, coaching, and savings support. Let \(q\in[0,1]\) be coverage, allowing implementation quality, costs, market prices, and spillovers to change with scale. Household \(i\)'s potential consumption at horizon \(T\) is \(C_{iT}(p,q)\). Define \(\tau_T(p,q)=E[C_{iT}(p,q)-C_{iT}(0,0)]\) for the entire target population, separately reporting beneficiaries and nonbeneficiaries. Let \(K(p,q)\) be total fiscal and implementation cost, and let \(B\) be an explicit budget. Characterize combinations maximizing a declared poverty-reduction functional subject to \(K(p,q)\le B\), and estimate whether lower-cost variants preserve durable gains. A factorial design can identify component interactions locally; national-scale effects additionally require variation in coverage and delivery constraints. Define durability using prespecified consumption, assets, and employment measures after support ends. Include next-generation outcomes if the follow-up can identify them. The authors' agenda explicitly asks about essential components and government scalability; this quantitative objective and its constraints are an adapted formulation, not a published optimization theorem.

\par\textbf{Formulation and scope.} The cited passage establishes a research direction, not this exact mathematical statement. The notation, population, policy menu, assumptions, and estimands are an original adaptation for this catalogue; no universal unsolved theorem or source priority is claimed.

\par\textbf{Open-question provenance.} An explicit question or research agenda was verified at the source locator. This does not by itself certify that every later version remains unresolved \sref{30007653}{1}.

\par\textbf{Historical priority.} Earliest directly inspected qualifying passage in this audit; earlier origins were not established. A review or research-agenda statement does not imply original priority.

\subsection{Short English statement}
In the explicitly specified setting: Which parts of multifaceted graduation programs are essential for durable and intergenerational poverty reduction when governments deliver them at national scale? The mathematical statement above fixes the target, assumptions and scope of this catalogue formulation.

\subsection{Sources}
\begin{itemize}\raggedright
\item[\textnormal{[S1]}]\hypertarget{source:30007653:1}{} Innovations for Poverty Action. 2025. The Ultra-Poor Graduation Approach: Research and Learning Agenda for the Next Wave of Graduation Programs. Section "What’s Next for Graduation Evidence", paragraphs 2–5.
\url{https://poverty-action.org/graduation-approach}.
{\small\textit{Earliest verified question/agenda reference; historical priority qualified below. The agenda explicitly asks which components are necessary, which adaptations reduce costs, and how government delivery scales.}}
\end{itemize}
\end{document}
